\section{Latent probe: does training teach the structure-only network the electronic structure?}
\label{si:probe}

If learning formation energies from structure alone made the network learn the electronic structure of the host,
a linear map from its learned description of the vacancy site and its host to the DFT descriptors should improve as
the network sees more training oxides, and should do better than the same map applied to an untrained network
reading the same geometry. This is claim C4, registered on the originally registered network (v1;
\texttt{docs/a4\_notes.md}). For every resample and budget, the readout input of the seed-0 S checkpoint (the
vacancy-site vector, the pooled host vector and the global state) was mapped by a ridge regression, fitted on the
training sites with its regularisation chosen by cross-validation grouped by host, to each descriptor given to D,
and scored by the coefficient of determination \Rsq{} on the test sites. The class means of the host-electronic and
the site-electronic descriptors are the quantities of the rule.

The registered rule (\texttt{docs/deviations.md}, the latent-probe row of step A4, committed before any v1 probe
number existed) has two conditions for each class. The trend condition asks that the recovered fraction of
descriptor variance rise with the number of training hosts: the Spearman correlation between the logarithm of the
budget and the class-mean \Rsq{} over all \numProbeCFourVOneRuleHostTrendNCells{} resample-budget cells must have an
interval above zero. The control condition asks that, at every budget from \numCfgProbePassRuleBeatsControlsMinBudget{}
hosts on, the trained probe exceed both the same probe fitted to descriptor rows permuted among the training
hosts (shuffled labels) and the probe on an untrained network of the same architecture (the registered
random-initialisation control), each difference with an interval above zero.

The trend condition holds for both classes, with Spearman correlations of \numProbeCFourVOneRuleHostTrendRho{}
\ci{\numProbeCFourVOneRuleHostTrendCiNinetyFiveItemOne{}}{\numProbeCFourVOneRuleHostTrendCiNinetyFiveItemTwo{}}
(host-electronic) and \numProbeCFourVOneRuleSiteTrendRho{}
\ci{\numProbeCFourVOneRuleSiteTrendCiNinetyFiveItemOne{}}{\numProbeCFourVOneRuleSiteTrendCiNinetyFiveItemTwo{}}
(site-electronic). The trained probe exceeds the shuffled-label control at every budget the rule names, by
\numProbeCFourVOneRuleHostControlsSixHundredFiftyFourShuffledMean{}
\ci{\numProbeCFourVOneRuleHostControlsSixHundredFiftyFourShuffledCiNinetyFiveItemOne{}}{\numProbeCFourVOneRuleHostControlsSixHundredFiftyFourShuffledCiNinetyFiveItemTwo{}}
for the host class at \numCfgBudgetsHostsItemSix{} hosts. It does not exceed the untrained network at any of them.
At \numCfgBudgetsHostsItemFour{} hosts it lies below it (host-electronic
\numProbeCFourVOneRuleHostControlsTwoHundredInitZeroMean{}
\ci{\numProbeCFourVOneRuleHostControlsTwoHundredInitZeroCiNinetyFiveItemOne{}}{\numProbeCFourVOneRuleHostControlsTwoHundredInitZeroCiNinetyFiveItemTwo{}},
site-electronic \numProbeCFourVOneRuleSiteControlsTwoHundredInitZeroMean{}
\ci{\numProbeCFourVOneRuleSiteControlsTwoHundredInitZeroCiNinetyFiveItemOne{}}{\numProbeCFourVOneRuleSiteControlsTwoHundredInitZeroCiNinetyFiveItemTwo{}}),
and at \numCfgBudgetsHostsItemSix{} hosts the differences are
\numProbeCFourVOneRuleHostControlsSixHundredFiftyFourInitZeroMean{}
\ci{\numProbeCFourVOneRuleHostControlsSixHundredFiftyFourInitZeroCiNinetyFiveItemOne{}}{\numProbeCFourVOneRuleHostControlsSixHundredFiftyFourInitZeroCiNinetyFiveItemTwo{}}
and \numProbeCFourVOneRuleSiteControlsSixHundredFiftyFourInitZeroMean{}
\ci{\numProbeCFourVOneRuleSiteControlsSixHundredFiftyFourInitZeroCiNinetyFiveItemOne{}}{\numProbeCFourVOneRuleSiteControlsSixHundredFiftyFourInitZeroCiNinetyFiveItemTwo{}},
intervals that include zero. C4 therefore does not hold, and under the registered rule the probe is reported here
rather than in the article (Figure~\ref{fig:probe}). The collaborator's probe on the primary network (v2, post hoc,
descriptive) shows the same pattern, with trend correlations of \numProbeDescriptiveSVTwoHostRuleTrendRho{} and
\numProbeDescriptiveSVTwoSiteRuleTrendRho{}, and C4 would not hold on v2 either. What is not established is that
training on formation energies adds linearly recoverable electronic information beyond what an untrained encoder of
the same structure already exposes. One reading, not tested here, is that the host-electronic descriptors are
largely functions of composition and structure, which an untrained encoder with element embeddings and distance
expansions already represents.

A sanity check confirms that the probe reads the vector that the network's final step reads. Applying the stored
readout of the checkpoint to the stored vectors reproduces the network's own test predictions to within
\numProbeFitVOneSanityZeroSixHundredFiftyFourHeadReplayMaxAbsDiffEV{}~eV at \numCfgBudgetsHostsItemSix{} hosts, the
level of single-precision arithmetic. A linear map from the same vectors to the network's own prediction, however,
recovers only part of it out of fold (\Rsq{} = \numProbeFitVOneSanityZeroSixHundredFiftyFourProbeToOwnPredictionRTwoOutOfFold{}
at \numCfgBudgetsHostsItemSix{} hosts, against \numProbeFitVOneSanityZeroTwentyFiveProbeToOwnPredictionRTwoOutOfFold{}
at \numCfgBudgetsHostsItemOne{}), because the readout is a nonlinear multilayer network and a linear probe need not
reproduce it. The probe therefore measures linear recoverability, a lower bound on what the vector encodes; this was
recorded as found and not adjusted (\texttt{results/probe\_fit\_v1.json}, key \texttt{payload.sanity}).

\begin{figure}[htbp]
\centering
\sifig{fig6_probe}
\caption{Latent probe: linear recovery of the DFT electronic descriptors from the learned description of the vacancy
site and its host in the structure-only network (registered analysis on the originally registered network, v1). A
linear map, fitted on the training hosts of each resample and budget, asks how well the learned vectors that the
network's final step reads (vacancy site, pooled host, global state) predict each descriptor given to D, scored by
\Rsq{} on the test sites. Trained probe: dashed line, open markers and hatched band (dashed bars in (c)). Controls in
grey: the same probe with descriptor rows permuted among training hosts (shuffled labels; squares) and the probe on
an untrained network of the same architecture (the registered random-initialisation control; triangles); two further
untrained networks (light grey, dotted, no band) show the spread of that control and are descriptive.
(\textbf{a}) Host-electronic and (\textbf{b}) site-electronic class-mean test \Rsq{} against training hosts on a
logarithmic axis: means over the resamples with two-sided intervals of the resample bootstrap; the dotted line is
zero. (\textbf{c}) Test \Rsq{} per descriptor at the largest budget, host-electronic descriptors above the grey rule
and site-electronic below. The registered rule does not hold, so the probe is placed in the Supplementary Materials.
The collaborator's probe on the primary network (v2) is post hoc and is not drawn. Records:
\texttt{results/probe\_c4\_v1.json}, \texttt{results/probe\_descriptive.json}.}
\label{fig:probe}
\end{figure}
\FloatBarrier
